\documentclass{ximera}
\title{Overzicht van notaties: vectoren, ladingen en veldlijnen}
\author{Daan Lipkens}

\begin{document}
\begin{abstract}
    Grafische vectorconventies gebruikt in de cursus.
\end{abstract}
\maketitle

Doorheen deze cursus gebruiken we telkens dezelfde kleuren en stijlen om vectoren, ladingen en veldlijnen visueel van elkaar te onderscheiden. Onderstaande tabel geeft een overzicht van deze notaties.

% ----------------------------------------------------------
% Hulpcommando's voor de voorbeeldjes in de tabel
\newcommand{\vecsample}[1]{%
    \begin{tikzpicture}[baseline=-0.5ex]
        \draw[#1] (0,0) -- (1.2,0);
    \end{tikzpicture}%
}
\newcommand{\lijnsample}[1]{%
    \begin{tikzpicture}[baseline=-0.5ex]
        \draw[#1] (0,0) -- (1.2,0);
    \end{tikzpicture}%
}
\newcommand{\ladingsample}[2]{%
    \begin{tikzpicture}[baseline=-0.6ex]
        \node[#1] at (0,0) {$#2$};
    \end{tikzpicture}%
}
% ----------------------------------------------------------

\begin{table}[h!]
    \caption{Overzicht van de gebruikte vector- en symboolnotaties.}
    \label{tab:notatie_overzicht}
    \centering
    \renewcommand{\arraystretch}{1.4}
    \begin{tabular}{@{}llll@{}}
        \toprule
        \multicolumn{4}{@{}l}{\textbf{Mechanica}} \\
        \midrule
        Positievector ($\vec{r}$) & \vecsample{positie} & Krachtvector ($\vec{F}$) & \vecsample{kracht} \\
        Snelheidsvector ($\vec{v}$) & \vecsample{snelheid} & Krachtcomponent ($\vec{F}_x$) & \vecsample{krachtcomponent} \\
        Snelheidscomponent ($\vec{v}_x$) & \vecsample{snelheidcomponent} & Resulterende kracht ($\vec{F}_\text{res}$) & \vecsample{resulterendekracht} \\
        Versnellingsvector ($\vec{a}$) & \vecsample{versnelling} & &\\
        Versnellingscomponent ($\vec{a}_x$) & \vecsample{versnellingcomponent} & & \\
        \midrule
        \multicolumn{4}{@{}l}{\textbf{Elektriciteit en magnetisme}} \\
        \midrule
        Positieve lading & \ladingsample{poslading}{+} & Negatieve lading & \ladingsample{neglading}{-} \\[0.5ex]
        Elektrisch veldvector ($\vec{E}$) & \vecsample{Evector} & Magnetisch veldvector ($\vec{B}$) & \vecsample{Mvector} \\
        Elektrisch veldcomponent ($\vec{E}_x$) & \vecsample{Evectorcomponent} & Magnetisch veldcomponent ($\vec{B}_x$) & \vecsample{Mvectorcomponent} \\
        Elektrische veldlijn ($E$) & \lijnsample{Eveldlijn} & Magnetische veldlijn ($B$) & \lijnsample{Mveldlijn} \\
        \bottomrule
    \end{tabular}
\end{table}

\end{document}